\renewcommand{\baselinestretch}{1.5}
\chapter*{\centering \large RESUMEN} % No numerado, centrado
\addcontentsline{toc}{chapter}{RESUMEN} 
\markboth{RESUMEN}{RESUMEN} % encabezado

El continuo despliegue de plantas solares hacia terrenos de topografía compleja genera alta variabilidad espacial de irradiancia, velocidad de viento, humedad y temperatura. Actualmente, la instrumentación basada en estaciones meteorológicas fijas asume un terreno plano y es incapaz de capturar esta diversidad, produciendo una alta incertidumbre en la estimación de la generación y el rendimiento energético de las instalaciones.
Frente a esta problemática, esta tesis propone el diseño conceptual de un robot autónomo terrestre para medir la variabilidad climática en plantas solares de topografía compleja. Esta solución dinámica reemplaza la necesidad de redes de sensores fijos, midiendo variables climáticas en múltiples puntos de manera focalizada para mejorar la evaluación energética en cumplimiento con la normativa internacional vigente.
Aplicando la metodología de diseño mecatrónico VDI 2221, se evaluaron diversas alternativas tecnológicas. El concepto óptimo validado es un vehículo robusto con tracción en sus cuatro ruedas, diseñado para mantener la estabilidad sobre terrenos irregulares. Para garantizar mediciones precisas, el diseño integra mecanismos articulados que elevan y orientan el sensor de irradiancia de manera independiente.
La navegación autónoma es gobernada por una unidad de procesamiento central que fusiona los datos del entorno mediante sensores de ubicación y movimiento. Utilizando algoritmos de control de trayectoria, el sistema logra seguir rutas planificadas y evadir obstáculos dinámicamente. Finalmente, la transmisión de los datos recolectados se realiza a través de una red de comunicación inalámbrica de largo alcance hacia el centro de control. En conclusión, el diseño propuesto demuestra ser una solución viable, descentralizada y robusta para las auditorías energéticas modernas.